
This work addressed the documentation adoption challenge by investigating whether AI-powered assistance could reduce completion friction. Our pilot with 75 users achieved 92\% median acceptance, substantially exceeding our conservative 70\% threshold and establishing that researchers will engage with AI-generated documentation suggestions.

Our contributions are: (1) Demonstration that researchers engage with AI documentation assistance at high rates (92\% acceptance with 22-percentage-point margin above threshold), validating the necessary precondition for pursuing quality improvements. (2) Evidence of robust cross-domain generalization without specialized tuning (1.0 percentage point variance across vision, NLP, and tabular datasets). (3) Validation that high acceptance reflects thoughtful engagement rather than passive acceptance (26.8\% modification rate).

This proof-of-concept establishes that researchers will use the tool but does not measure whether use improves documentation outcomes. Full-scale deployment with A/B testing is needed to measure whether 92\% acceptance translates to improved completeness, higher expert quality ratings, and reduced time-to-completion.

Future work includes measuring downstream documentation quality improvements through controlled deployment, understanding longitudinal acceptance patterns over 6-12 months, validating production corpus quality effects through systematic curation and ablation studies, extending to other research documentation tasks (IRB proposals, grant applications, methodology sections), assessing multilingual and cross-cultural generalization, and evaluating adversarial resistance with controlled studies using minimal-effort incentives.

What this work demonstrates is that application domain characteristics matter for AI assistance adoption. Documentation's lower error stakes and higher time pressure may create more favorable conditions for AI acceptance than code generation, where correctness requirements demand careful verification. This insight suggests AI assistance research should stratify applications by error tolerance rather than treating all content generation as equivalent.

The path forward involves validating the complete value chain from engagement to quality improvement. Our proof-of-concept establishes that the first step is achievable—researchers will use the tool—providing foundation for measuring whether use translates to better documentation.
